\begin{textsample}{POS Dim 3 – vem\_ed\_34 – Score 3.00 – vem\_ed\_34/t181.txt}  \label{ex:f3_pos_062}
Aos Leitores Osvaldo Succi Jr. - Coordenador de Projetos Este número de VEm convida a refletir sobre a consolidação e o futuro dos Projetos Colaborativos Internacionais (PCIs) desenvolvidos nas Fatecs.
Em março de 2026, representei o Centro Paula Souza (CPS) e a América Latina no webinar organizado pela Fundação COIL Connect, que celebrou os 20 anos da abordagem COIL (Collaborative Online International Learning).
O evento gratuito reuniu vozes de diferentes continentes para \textbf{discutir} conquistas, desafios e o papel central das relações humanas na educação global.
Na FAUBAI 2026, apresentei três trabalhos, evidenciando experiências que reafirmam o impacto dos PCIs na educação profissional e tecnológica.
A edição apresenta ainda um PCI realizado entre a Fatec Americana e o CEAM (Peru) na área têxtil, ilustrando como o Intercâmbio Virtual promove aprendizagem significativa, diálogo intercultural e desenvolvimento profissional.
Para produzir o Artigo de Opinião, reunimos a equipe de Apoio à Internacionalização do Ensino Superior da Divisão de Extensão e Pesquisa da Coordenadoria Geral de Ensino Superior de Graduação do CPS e dialogamos sobre aprendizagens acumuladas pela equipe, \textbf{apontando} fatores de \textbf{sucesso} dos PCIs e cuidados no planejamento, execução e avaliação, para que se tornem colaborações duradouras.
Deseja contribuir com Artigo de Opinião?
Escreva para cgesg.pci@cps.sp.gov.br.
Boa leitura!

% matched lemmas: apontar, discutir, sucesso
\end{textsample}
